\subsection{Discussion}
\label{sec:discussion}

Table~\ref{tab:main-results} reports \emph{relative} baseline comparison under the documented simulated latency model (\S\ref{sec:setup}); absolute RTT values are not live hospital traces.
B2 reduces M1 p50 from 374.0\,ms (B0) to 208.3\,ms and raises M2 from 0.560 to 0.604 on ShezhenV3 test frames with physical resampling and ROI Edge-IQA ($\tau_{B2}{=}0.465$).
All three seeds show B2 outperforming B0 on M1 and M2; low M2 variance reflects deterministic test-set scores given a shared frame plan, with remaining randomness from the latency model and frame-plan shuffling.

COCO tongue ROI (8\% padded union box) aligns scoring with downstream anatomy but does not eliminate clear/blur overlap on Laplacian features (Table~\ref{tab:roi-ablation}, negative separation $-0.541$).
Because ShezhenV3 lacks human blur labels, we inject Gaussian defocus on half of test frames, mirroring validation calibration; M5 Spearman $\rho{\approx}0.36$--$0.47$ reflects residual overlap under real texture.
B3 (MobileNetV3-Small, 98.3\% val accuracy) achieves M2$=0.934$ at $\tau_{B2}$ versus B2 M2$=0.604$; we retain B2 for lowest M1 p50 and auditable \texttt{flags}, while B3/B4 suit higher valid-rate targets with extra edge compute.

Edge-IQA and LangGraph are hardware-agnostic: any edge gateway exposing HTTP \texttt{/vision/evaluate} and skill endpoints can replay the same JSONL protocol.
M2 measures frames passing the scorer threshold---a routing proxy, not independent clinical quality annotation; future work should correlate M2 with downstream classifier confidence on accepted vs.\ rejected frames.
Extended tables (NR-IQA baselines, per-seed breakdown, deployment notes, failure modes) are provided in the Supplementary Material.

The results should be read as a deployment trade-off.
B0 is simple but spends cloud RTT on frames that could have been repaired locally; B1 saves some upload time but cannot improve frame validity because it lacks a physical recovery action.
B2 is the smallest closed-loop policy that changes both latency and input quality, and its gain on blur-injected frames (Table~\ref{tab:stratified-m2}) shows that improvement comes from resampling rather than from re-labeling the same images.
B3/B4 indicate how much valid-frame rate remains available if a site accepts learned-model maintenance and higher edge RTT.
This separation is important for medical robotics: the conservative B2 path is easier to certify and audit, while the learned tier can be enabled only when local validation justifies the added model risk.
